\documentclass[a4paper, 10pt, colorlinks=true, linkcolor=black, urlcolor=blue, citecolor=blue, hidelinks]{article}
\usepackage[table]{xcolor}
\usepackage[italian]{babel} 
\usepackage{comment}      
\usepackage{graphicx} 
\usepackage{capt-of}
\usepackage{multirow}
\usepackage{longtable}
\usepackage{geometry}
\usepackage{xurl}
\usepackage[nopatch=footnote]{microtype}
\usepackage{bookmark}
\usepackage{float} 
\usepackage{setspace}
\usepackage{tabularx}
\usepackage{fancyhdr}
\usepackage{array}

\definecolor{bluBitless}{HTML}{3A75C4}
\newcolumntype{L}{>{\raggedright\arraybackslash}X}

\pagestyle{fancy}
\fancyhf{}
\fancyhead[L]{The Bitless}
\fancyhead[R]{Piano di Qualifica}
\fancyfoot[C]{\thepage}

\usepackage{hyperref}
\hypersetup{
    colorlinks=true,
    linkcolor=black, 
    urlcolor=bluBitless,     
    citecolor=black,  
    filecolor=black,
    breaklinks=true,
    pdfpagemode=UseNone    
}

\begin{document}
\raggedbottom
\setlength{\emergencystretch}{2em}

\newcommand{\titolo}{Piano di Qualifica$_G$}
\newcommand{\dataverbale}{}
\newcommand{\dataversione}{2026-06-04}
\newcommand{\versione}{0.1.1}
\newcommand{\stato}{Stesura}
\newcommand{\destinatari}{The Bitless, prof. T. Vardanega, prof. R. Cardin, Zucchetti}
\def\templatedir{../../../templates}

\setcounter{secnumdepth}{5}
\setcounter{tocdepth}{5}

\begin{titlepage}
\pagenumbering{gobble}
\begin{titlepage} 
      \begin{center}
            \thispagestyle{empty}
            \providecommand{\templatedir}{../../../../templates}
            
\includegraphics[width=15cm]{\templatedir/logo_gruppo.png}
            \begin{spacing}{1.2}
                \Huge{\textbf{\titolo \ \dataverbale}}
            \end{spacing}

            \vspace{2cm}


            \begin{center}
            \textbf{The Bitless} (Gruppo 18): \\
            Lorenzo Battistella (2103116), Edoardo De Piccoli (2101055), Davide Facco (2087852), \\ Anna Marini (2110986), Andrea Menegaldo (2116426), Samuele Vendramin (2111934),\\ Simone Zecchinato (2113189)
            \end{center}

            \vspace{1cm}

            \begin{center}
                \begin{tabular}{|c|c|c|c|}
                    \hline
                    \multicolumn{4}{|c|}{\textbf{Informazioni documento}} \\
                    \hline
                    \textbf{Versione} & \textbf{Data} & \textbf{Stato}  & \textbf{Destinatari} \\
                    \hline
                    \versione & \dataversione & \stato & \destinatari \\
                    \hline
                \end{tabular}
            \end{center}

            \vfill
            
            \begin{center}
                Contatti: thebitless.swe@gmail.com
            \end{center}
            
        \end{center}
    \end{titlepage}
\end{titlepage}

\newpage
\pagenumbering{arabic}

\begingroup
\small
\renewcommand{\arraystretch}{1.25}
\begin{longtable}{|c|c|c|c|c|p{4cm}|}
    \hline
    \multicolumn{6}{|c|}{\textbf{Registro delle modifiche}} \\
    \hline
    \textbf{Versione} & \textbf{Data} & \textbf{Autore} & \textbf{Verificatore} & \textbf{Approvatore} & \textbf{Descrizione} \\
    \hline
    \endhead
    0.3.0 & 2026-06-08 & S. Zecchinato & - & - & Aggiunge valutazione MPC-AC/MPC-ETC \\ 
    \hline
    0.2.0 & 2026-06-07 & S. Zecchinato & - & - & Aggiunge valutazione MPC-PV/MPC-EV \\
    \hline
    0.1.1 & 2026-06-04 & L. Battistella & - & - & Aggiornamento test$_G$ di sistema$_G$. \\ 
    \hline
    0.1.0 & 2026-05-14 & A. Marini & - & - & Creazione prima versione del documento Piano di Qualifica$_G$. Inserimento struttura iniziale, riferimenti, metriche di qualità e strategia di testing. \\ \hline

\end{longtable}
\endgroup

\newpage

\tableofcontents
\newpage

\listoftables
\newpage

\listoffigures
\newpage

\section{Introduzione}

\subsection{Scopo del documento}

Il presente documento descrive il \textit{Piano di Qualifica} adottato dal gruppo \textit{The Bitless} per il progetto \textit{Second Brain}, relativo al capitolato$_G$ C6 proposto dall'azienda Zucchetti.

Il Piano di Qualifica$_G$ ha lo scopo di definire in modo esplicito gli obiettivi di qualità perseguiti durante lo sviluppo$_G$ del prodotto software e le modalità con cui tali obiettivi vengono misurati, controllati e progressivamente migliorati.

Il documento raccoglie quindi:
\begin{itemize}
    \item le metriche utilizzate per valutare la qualità dei processi interni al gruppo;
    \item le metriche utilizzate per valutare la qualità del prodotto software;
    \item la strategia di verifica$_G$ e validazione$_G$ adottata;
    \item la classificazione dei test$_G$ previsti;
    \item i criteri con cui interpretare i risultati raccolti nel cruscotto di valutazione.
\end{itemize}

Il Piano di Qualifica$_G$ non ha l'obiettivo di sostituire le \textit{Norme di Progetto}, nelle quali vengono definite le procedure operative seguite dal gruppo. Esso stabilisce invece quali aspetti devono essere misurati, quali valori sono considerati accettabili e quali valori rappresentano un risultato ottimale.

\subsection{Scopo del prodotto}

Il prodotto richiesto dal capitolato$_G$ C6 consiste in un'applicazione web per la scrittura e la gestione di note in formato Markdown, arricchita dall'integrazione$_G$ con un Large Language Model$_G$.

L'applicazione deve permettere all'utente di:
\begin{itemize}
    \item scrivere testo Markdown in un'area di editing;
    \item visualizzare il risultato formattato in una zona di rendering;
    \item applicare operazioni basate su LLM$_G$ al testo completo o a una selezione;
    \item riassumere, riscrivere e tradurre il testo;
    \item ottenere analisi critiche secondo il modello dei sei cappelli per pensare$_G$;
    \item generare un testo completo a partire da un prompt;
    \item salvare e recuperare note in formato testuale.
\end{itemize}

La qualità del prodotto deve quindi essere valutata non solo rispetto alla correttezza funzionale, ma anche rispetto alla facilità d'uso, alla stabilità dell'interazione con il modello linguistico, alla manutenibilità$_G$ del codice e alla verificabilità delle funzionalità principali.

\subsection{Glossario}

Per evitare ambiguità terminologiche, i termini tecnici, di dominio$_G$ o potenzialmente equivoci sono definiti nel documento \textit{Glossario}. Alla prima occorrenza significativa nel presente documento, tali termini devono essere interpretati secondo la definizione riportata nel Glossario.

\subsection{Riferimenti}

\subsubsection{Riferimenti normativi}

\begin{itemize}
    \item \textbf{Norme di Progetto v0.1.0} \\
    Documento interno del gruppo \textit{The Bitless}, contenente le procedure operative, gli strumenti adottati e le convenzioni seguite nello svolgimento del progetto.

    \item \textbf{Piano di Progetto v0.1.0} \\
    Documento interno del gruppo \textit{The Bitless}, contenente la pianificazione delle attività, la stima dei costi, la gestione dei rischi e il consuntivo di avanzamento.

    \item \textbf{Analisi dei Requisiti v2.0.0} \\
    Documento esterno del gruppo \textit{The Bitless}, contenente la descrizione dei casi d'uso$_G$, dei Requisiti Funzionali$_G$, dei requisiti$_G$ non funzionali e dei vincoli di progetto.

    \item \textbf{Capitolato C6 -- Second Brain} \\
    Capitolato$_G$ proposto dall'azienda Zucchetti per il corso di Ingegneria del Software 2025/2026.
\end{itemize}

\subsubsection{Riferimenti informativi}

\begin{itemize}
    \item \textbf{Glossario v0.1.0} \\
    Documento contenente le definizioni dei termini rilevanti utilizzati nei prodotti documentali del gruppo.

    \item \textbf{Slide T07 -- Qualità del software} \\
    Materiale didattico del corso di Ingegneria del Software.

    \item \textbf{Slide T08 -- Qualità del software} \\
    Materiale didattico del corso di Ingegneria del Software.

    \item \textbf{ISO/IEC 25010} \\
    Standard utilizzato come riferimento informativo per la classificazione delle caratteristiche di qualità del prodotto software.

    \item \textbf{Documentazione Markdown} \\
    Riferimento informativo per la sintassi Markdown e per le funzionalità di rendering previste dal prodotto.

    \item \textbf{Documentazione API compatibili con OpenAI} \\
    Riferimento informativo per l'integrazione$_G$ con modelli linguistici tramite chiamate HTTP.
\end{itemize}

\section{Metriche di qualità del processo}

Le metriche di qualità del processo permettono di controllare il modo in cui il gruppo organizza e svolge le proprie attività. Esse non misurano direttamente il comportamento del prodotto finale, ma valutano l'efficacia della pianificazione, la stabilità dei requisiti$_G$, la qualità della documentazione, l'efficienza delle verifiche e la sostenibilità del lavoro.

Nel presente documento le metriche di processo sono identificate con il prefisso \textbf{MPC}, acronimo di \textit{Metrica di Processo e Controllo}. Ogni metrica è associata a:
\begin{itemize}
    \item un identificativo univoco;
    \item un nome descrittivo;
    \item una soglia di accettabilità;
    \item un valore ottimale.
\end{itemize}

Le soglie riportate hanno lo scopo di aiutare il gruppo a interpretare i dati raccolti durante il progetto. Il raggiungimento del valore accettabile indica che il processo è sotto controllo; il raggiungimento del valore ottimale indica invece una condizione pienamente soddisfacente.

\subsection{Fornitura}

Le metriche relative alla fornitura misurano l'andamento del progetto rispetto alla pianificazione economica e temporale. Esse permettono di confrontare quanto era stato previsto con quanto è stato effettivamente realizzato.

\begingroup
\small
\renewcommand{\arraystretch}{1.25}
\begin{longtable}{|c|p{4cm}|p{3cm}|p{3cm}|}
    \hline
    \textbf{ID} & \textbf{Nome metrica} & \textbf{Valore accettabile} & \textbf{Valore ottimale} \\
    \hline
    \endhead

    MPC-PV & Planned Value & $> 0$ & Coerente con la pianificazione approvata \\ \hline
    MPC-EV & Earned Value & $\geq 80\%$ del valore pianificato & $\geq$ Planned Value \\ \hline
    MPC-AC & Actual Cost$_G$ & $\leq 120\%$ del costo pianificato & $\leq$ costo pianificato \\ \hline
    MPC-CV & Cost Variance & $\geq -10\%$ & $\geq 0$ \\ \hline
    MPC-SV & Schedule Variance & $\geq -10\%$ & $\geq 0$ \\ \hline
    MPC-EAC & Estimate At Completion & $\leq 110\%$ del budget$_G$ & $\leq$ budget$_G$ approvato \\ \hline
    MPC-ETC & Estimate To Complete & Valore definito e aggiornato & Valore coerente con il piano a finire \\ \hline
    MPC-TCR & Task$_G$ Completion Rate & $\geq 85\%$ & $100\%$ \\ \hline
    MPC-TS & Task$_G$ Slippage & $\leq 15\%$ & $0\%$ \\ \hline

    \caption{Metriche per la qualità del processo di fornitura.}
    \label{tab:metriche-fornitura}
\end{longtable}
\endgroup

\subsection{Sviluppo}

Le metriche relative allo sviluppo$_G$ controllano la stabilità del lavoro tecnico e la capacità del gruppo di introdurre modifiche in modo ordinato, verificabile e tracciabile.

Nel progetto \textit{Second Brain}, particolare attenzione viene posta alla stabilità dei requisiti$_G$ collegati all'integrazione$_G$ con LLM$_G$, poiché la natura esplorativa del capitolato$_G$ può portare a raffinamenti progressivi dei prompt, delle funzioni disponibili e delle modalità di interazione con il testo.

\begingroup
\small
\renewcommand{\arraystretch}{1.25}
\begin{longtable}{|c|p{4cm}|p{3cm}|p{3cm}|}
    \hline
    \textbf{ID} & \textbf{Nome metrica} & \textbf{Valore accettabile} & \textbf{Valore ottimale} \\
    \hline
    \endhead

    MPC-RSI & Requirements Stability Index & $\geq 80\%$ & $100\%$ \\ \hline
    MPC-PRCT & Pull Request$_G$ Cycle Time & $\leq 72$ ore & $\leq 24$ ore \\ \hline
    MPC-MR & Merge$_G$ Rework & $\leq 20\%$ delle pull request$_G$ & $\leq 5\%$ delle pull request$_G$ \\ \hline
    MPC-IR & Issue$_G$ Reopening Rate & $\leq 15\%$ & $0\%$ \\ \hline

    \caption{Metriche per la qualità del processo di sviluppo.}
    \label{tab:metriche-sviluppo}
\end{longtable}
\endgroup

\subsection{Documentazione}

Le metriche relative alla documentazione misurano la leggibilità, la correttezza e la verificabilità dei documenti prodotti. Poiché i documenti rappresentano il principale mezzo di comunicazione verso committente$_G$, proponenti e docenti, essi devono essere chiari, aggiornati e coerenti tra loro.

\begingroup
\small
\renewcommand{\arraystretch}{1.25}
\begin{longtable}{|c|p{4cm}|p{3cm}|p{3cm}|}
    \hline
    \textbf{ID} & \textbf{Nome metrica} & \textbf{Valore accettabile} & \textbf{Valore ottimale} \\
    \hline
    \endhead

    MPC-IG & Indice di Gulpease & $\geq 40$ & $\geq 60$ \\ \hline
    MPC-CO & Correttezza ortografica & Nessun errore grave & $0$ errori \\ \hline
    MPC-CD & Completezza documentale & $\geq 90\%$ sezioni previste & $100\%$ sezioni previste \\ \hline
    MPC-CR & Coerenza dei riferimenti & $\geq 95\%$ riferimenti validi & $100\%$ riferimenti validi \\ \hline

    \caption{Metriche per la qualità della documentazione.}
    \label{tab:metriche-documentazione}
\end{longtable}
\endgroup

\subsection{Verifica}

Le metriche relative alla verifica$_G$ misurano l'efficacia delle attività di controllo svolte sul prodotto e sui documenti. Esse permettono di valutare se gli errori vengono individuati tempestivamente e se le correzioni introdotte risultano effettivamente risolutive.

\begingroup
\small
\renewcommand{\arraystretch}{1.25}
\begin{longtable}{|c|p{4cm}|p{3cm}|p{3cm}|}
    \hline
    \textbf{ID} & \textbf{Nome metrica} & \textbf{Valore accettabile} & \textbf{Valore ottimale} \\
    \hline
    \endhead

    MPC-CC & Code Coverage & $\geq 70\%$ & $\geq 90\%$ \\ \hline
    MPC-TSR & Test$_G$ Success Rate & $\geq 95\%$ & $100\%$ \\ \hline
    MPC-DD & Defect Density & $\leq 2$ difetti per componente & $0$ difetti aperti \\ \hline
    MPC-TR & Test$_G$ Requirements Coverage & $100\%$ requisiti$_G$ obbligatori coperti & $100\%$ requisiti$_G$ coperti \\ \hline
    MPC-QMS & Quality Metrics Satisfied & $\geq 80\%$ & $100\%$ \\ \hline

    \caption{Metriche per la qualità del processo di verifica.}
    \label{tab:metriche-verifica}
\end{longtable}
\endgroup

\section{Metriche di qualità del prodotto}

Le metriche di qualità del prodotto misurano le caratteristiche del software realizzato. A differenza delle metriche di processo, esse valutano direttamente il comportamento dell'applicazione, la sua affidabilità, la sua facilità d'uso e la sua capacità di soddisfare i requisiti$_G$ del capitolato$_G$.

Nel presente documento le metriche di prodotto sono identificate con il prefisso \textbf{MPD}, acronimo di \textit{Metrica di Prodotto}. Le caratteristiche considerate sono:
\begin{itemize}
    \item funzionalità;
    \item affidabilità;
    \item usabilità;
    \item efficienza;
    \item manutenibilità$_G$;
    \item sicurezza.
\end{itemize}

\subsection{Funzionalità}

Le metriche di funzionalità misurano quanto il prodotto soddisfa i requisiti$_G$ individuati nell'Analisi dei Requisiti$_G$. Per il progetto \textit{Second Brain}, la copertura dei requisiti$_G$ obbligatori è particolarmente importante, poiché il capitolato$_G$ richiede un insieme minimo di funzionalità necessarie per considerare accettabile il prodotto.

\begingroup
\small
\renewcommand{\arraystretch}{1.25}
\begin{longtable}{|c|p{4cm}|p{3cm}|p{3cm}|}
    \hline
    \textbf{ID} & \textbf{Nome metrica} & \textbf{Valore accettabile} & \textbf{Valore ottimale} \\
    \hline
    \endhead

    MPD-CRO & Copertura requisiti$_G$ obbligatori & $100\%$ & $100\%$ \\ \hline
    MPD-CRD & Copertura requisiti$_G$ desiderabili & $\geq 50\%$ & $\geq 80\%$ \\ \hline
    MPD-CROp & Copertura requisiti$_G$ opzionali & $\geq 0\%$ & $\geq 60\%$ \\ \hline
    MPD-CLLM & Copertura funzioni LLM$_G$ obbligatorie & $100\%$ & $100\%$ \\ \hline
    MPD-CMD & Correttezza rendering Markdown & $\geq 95\%$ casi previsti & $100\%$ casi previsti \\ \hline
    MPD-SN & Salvataggio e recupero note & $100\%$ casi obbligatori & $100\%$ casi obbligatori \\ \hline

    \caption{Metriche per la qualità funzionale del prodotto.}
    \label{tab:metriche-funzionalita}
\end{longtable}
\endgroup

\subsection{Affidabilità}

Le metriche di affidabilità misurano la capacità del prodotto di mantenere un comportamento corretto anche in presenza di input$_G$ non ottimali, errori di comunicazione o risposte inattese da parte del modello linguistico.

Nel contesto del capitolato$_G$ C6, l'affidabilità riguarda in particolare:
\begin{itemize}
    \item la stabilità dell'editor durante la scrittura;
    \item la conservazione del testo inserito dall'utente;
    \item la gestione degli errori provenienti dall'LLM$_G$;
    \item la capacità di evitare perdita di dati durante salvataggio e caricamento delle note.
\end{itemize}

\begingroup
\small
\renewcommand{\arraystretch}{1.25}
\begin{longtable}{|c|p{4cm}|p{3cm}|p{3cm}|}
    \hline
    \textbf{ID} & \textbf{Nome metrica} & \textbf{Valore accettabile} & \textbf{Valore ottimale} \\
    \hline
    \endhead

    MPD-DL & Data Loss & $0$ perdite nei casi principali & $0$ perdite in ogni test$_G$ \\ \hline
    MPD-EH & Error Handling & $\geq 90\%$ errori gestiti & $100\%$ errori gestiti \\ \hline
    MPD-SR & Session Recovery & Recupero testo nei casi previsti & Recupero completo e automatico \\ \hline
    MPD-AR & Availability Rate & $\geq 95\%$ durante i test$_G$ & $\geq 99\%$ durante i test$_G$ \\ \hline

    \caption{Metriche per la qualità dell'affidabilità.}
    \label{tab:metriche-affidabilita}
\end{longtable}
\endgroup

\subsection{Usabilità}

Le metriche di usabilità valutano la facilità con cui l'utente può comprendere e utilizzare l'applicazione. Il prodotto deve risultare accessibile anche a utenti non esperti, come richiesto dalla natura del capitolato$_G$.

L'interfaccia deve permettere di distinguere chiaramente:
\begin{itemize}
    \item l'area di editing Markdown;
    \item l'area di rendering;
    \item i comandi base sul testo;
    \item i comandi collegati all'LLM$_G$;
    \item le funzioni di salvataggio e caricamento.
\end{itemize}

\begingroup
\small
\renewcommand{\arraystretch}{1.25}
\begin{longtable}{|c|p{4cm}|p{3cm}|p{3cm}|}
    \hline
    \textbf{ID} & \textbf{Nome metrica} & \textbf{Valore accettabile} & \textbf{Valore ottimale} \\
    \hline
    \endhead

    MPD-LT & Learning Time & $\leq 30$ minuti & $\leq 10$ minuti \\ \hline
    MPD-NE & Navigation Errors & $\leq 3$ errori per scenario & $0$ errori per scenario \\ \hline
    MPD-UI & UI Consistency & $\geq 90\%$ elementi coerenti & $100\%$ elementi coerenti \\ \hline
    MPD-FU & Feature Understandability & $\geq 80\%$ funzioni comprese senza guida & $\geq 95\%$ funzioni comprese senza guida \\ \hline

    \caption{Metriche per la qualità dell'usabilità.}
    \label{tab:metriche-usabilita}
\end{longtable}
\endgroup

\subsection{Efficienza}

Le metriche di efficienza misurano i tempi di risposta e l'uso delle risorse. Nel progetto \textit{Second Brain}, una parte significativa del tempo di risposta può dipendere dal modello linguistico utilizzato; per questo motivo è opportuno distinguere tra:
\begin{itemize}
    \item tempo di risposta dell'interfaccia locale;
    \item tempo di rendering Markdown;
    \item tempo di risposta delle chiamate LLM$_G$.
\end{itemize}

\begingroup
\small
\renewcommand{\arraystretch}{1.25}
\begin{longtable}{|c|p{4cm}|p{3cm}|p{3cm}|}
    \hline
    \textbf{ID} & \textbf{Nome metrica} & \textbf{Valore accettabile} & \textbf{Valore ottimale} \\
    \hline
    \endhead

    MPD-RT & Response Time interfaccia & $\leq 2$ secondi & $\leq 0{,}5$ secondi \\ \hline
    MPD-MRT & Markdown Rendering Time & $\leq 1$ secondo & $\leq 0{,}3$ secondi \\ \hline
    MPD-LLMRT & LLM$_G$ Response Time & $\leq 15$ secondi & $\leq 8$ secondi \\ \hline
    MPD-LS & Loading Speed & $\leq 3$ secondi & $\leq 1$ secondo \\ \hline

    \caption{Metriche per la qualità dell'efficienza.}
    \label{tab:metriche-efficienza}
\end{longtable}
\endgroup

\subsection{Manutenibilità}

Le metriche di manutenibilità$_G$ valutano quanto il codice sia comprensibile, modificabile e testabile. Considerata la natura evolutiva del progetto, è importante che l'architettura permetta di aggiungere o modificare funzioni LLM$_G$, prompt e modalità di salvataggio senza introdurre regressioni.

\begingroup
\small
\renewcommand{\arraystretch}{1.25}
\begin{longtable}{|c|p{4cm}|p{3cm}|p{3cm}|}
    \hline
    \textbf{ID} & \textbf{Nome metrica} & \textbf{Valore accettabile} & \textbf{Valore ottimale} \\
    \hline
    \endhead

    MPD-CYC & Complessità ciclomatica & $\leq 15$ & $\leq 10$ \\ \hline
    MPD-CS & Code Smell & $\leq 5$ per componente & $\leq 1$ per componente \\ \hline
    MPD-DUP & Duplicazione del codice & $\leq 10\%$ & $\leq 3\%$ \\ \hline
    MPD-MOD & Modularità & Componenti principali separate & Dipendenze ridotte e ben definite \\ \hline

    \caption{Metriche per la qualità della manutenibilità.}
    \label{tab:metriche-manutenibilita}
\end{longtable}
\endgroup

\subsection{Sicurezza}

Le metriche di sicurezza misurano la capacità del prodotto di proteggere dati, chiavi di accesso e contenuti dell'utente. Poiché l'applicazione interagisce con servizi LLM$_G$ esterni, è necessario evitare l'esposizione impropria di credenziali e gestire con attenzione i dati inviati tramite API$_G$.

\begingroup
\small
\renewcommand{\arraystretch}{1.25}
\begin{longtable}{|c|p{4cm}|p{3cm}|p{3cm}|}
    \hline
    \textbf{ID} & \textbf{Nome metrica} & \textbf{Valore accettabile} & \textbf{Valore ottimale} \\
    \hline
    \endhead

    MPD-SK & Secret Key Exposure & $0$ chiavi esposte & $0$ chiavi esposte \\ \hline
    MPD-IV & Input$_G$ Validation & $\geq 95\%$ input$_G$ validati & $100\%$ input$_G$ validati \\ \hline
    MPD-DS & Data Sanitization & $\geq 95\%$ dati sanificati & $100\%$ dati sanificati \\ \hline
    MPD-PR$_G$ & Privacy Risk & Nessun dato sensibile non necessario inviato & Minimizzazione completa dei dati inviati \\ \hline

    \caption{Metriche per la qualità della sicurezza.}
    \label{tab:metriche-sicurezza}
\end{longtable}
\endgroup

\section{Strategia di verifica e testing}

La strategia di verifica$_G$ adottata dal gruppo ha lo scopo di fornire evidenza oggettiva del corretto avanzamento del progetto e della conformità del prodotto rispetto ai requisiti$_G$ individuati.

Le attività di testing vengono organizzate in modo progressivo. Nelle fasi iniziali vengono privilegiati test$_G$ e controlli sulle funzionalità fondamentali, mentre nelle fasi successive la copertura viene estesa alle interazioni tra componenti, ai flussi completi e ai requisiti$_G$ non funzionali.

Ogni test$_G$ è identificato da un codice univoco, da una descrizione, dal requisito associato e da uno stato.

\subsection{Stati dei test}

Gli stati utilizzati per classificare i test$_G$ sono i seguenti:

\begin{itemize}
    \item \textbf{NI -- Non Implementato}: il test$_G$ è stato previsto, ma non è ancora stato implementato;
    \item \textbf{I -- Implementato}: il test$_G$ è stato implementato, ma non ancora eseguito o verificato stabilmente;
    \item \textbf{S -- Superato}: il test$_G$ è stato eseguito e ha prodotto esito positivo;
    \item \textbf{NS -- Non Superato}: il test$_G$ è stato eseguito, ma ha prodotto esito negativo;
    \item \textbf{NA -- Non Applicabile}: il test$_G$ non è applicabile alla versione corrente del prodotto.
\end{itemize}

\subsection{Copertura del codice}

La copertura del codice misura la percentuale di codice sorgente eseguita durante l'esecuzione dei test$_G$ automatici.

Questa metrica viene utilizzata come indicatore della quantità di codice effettivamente verificata, ma non è sufficiente da sola a garantire la correttezza del prodotto. Per questo motivo viene interpretata insieme al tasso di successo dei test$_G$, alla copertura dei requisiti$_G$ e al numero di difetti individuati.

Il valore minimo accettabile è fissato al $70\%$, mentre il valore ottimale è pari o superiore al $90\%$.

\subsection{Test unitari}

I Test Unitari$_G$ verificano singole unità software$_G$ in isolamento. Nel progetto \textit{Second Brain}, tali test$_G$ riguardano principalmente:
\begin{itemize}
    \item parsing e gestione del testo Markdown;
    \item funzioni di aggiornamento dello stato dell'editor;
    \item validazione$_G$ degli input$_G$;
    \item costruzione dei prompt da inviare all'LLM$_G$;
    \item funzioni di salvataggio e lettura delle note;
    \item gestione degli errori applicativi.
\end{itemize}

La specifica dei Test Unitari$_G$ viene raffinata con l'avanzamento della progettazione di dettaglio. In questa fase, il Piano di Qualifica$_G$ definisce la strategia e gli obiettivi di copertura, mentre l'elenco puntuale dei test$_G$ sarà aggiornato quando le unità software$_G$ saranno stabilizzate.

\subsection{Test di integrazione}

\newcolumntype{D}{>{\hsize=1.45\hsize}X} % Descrizione Interfaccia
\newcolumntype{M}{>{\hsize=0.55\hsize}X} % Moduli Coinvolti

I test$_G$ di integrazione$_G$ del PoC$_G$ verificano che i moduli dell'architettura dialoghino correttamente ai loro confini, validando interfacce, formato dei dati scambiati e propagazione degli errori lungo l'intera catena.
\begin{table}[H]
	\centering
	\begin{tabularx}{\textwidth}{|l|D|M|}
		\hline
		\textbf{Codice} & \textbf{Descrizione Interfaccia} & \textbf{Moduli Coinvolti} \\
		\hline
		\textbf{TI-001} & Verifica$_G$ dello scambio dati e della gestione degli errori tra hook e FastAPI; invio della richiesta via fetch POST, ricezione dello stream SSE e gestione dell'indisponibilità del provider (HTTP 503). & Frontend$_G$ React (hook useLlmStream) $\leftrightarrow$ Backend FastAPI \\
		\hline
		\textbf{TI-002} & Verifica$_G$ dell'invio del prompt di sistema$_G$ e del contesto utente verso il modello LLM$_G$ e della ricezione dello stream di risposta openai-compatible dall'API$_G$ del provider, tramite client HTTP asincrono (httpx). & LLMClient / LiteLLMClient $\leftrightarrow$ Provider LiteLLM Zucchetti \\
		\hline
		\textbf{TI-003} & Verifica$_G$ della propagazione dell'annullamento lungo l'intera catena; dall'AbortController lato client a Request.is\_disconnected() lato server, fino alla chiusura dello stream del provider. & AbortController (client) $\leftrightarrow$ Request.is\_disconnected() (server) \\
		\hline
		\textbf{TI-004} & Verifica$_G$ del parsing dei chunk SSE e dell'aggiornamento progressivo dello stato globale dell'applicazione (appendChunk / streamedOutput). & useLlmStream (parsing via lib/sse.ts) $\leftrightarrow$ Store Zustand \\
		\hline
		\textbf{TI-005} & Verifica$_G$ del binding del testo e del re-render granulare dell'anteprima durante lo streaming, mediato dalla store globale: l'editor mantiene il focus mentre solo la preview si aggiorna. & Store Zustand (selettori granulari) $\leftrightarrow$ Editor / Preview \\
		\hline
		\textbf{TI-006} & Verifica$_G$ dell'applicazione a runtime della politica CORS tra i due servizi containerizzati; le richieste provenienti da un'origine non presente nella whitelist vengono bloccate, quelle da origine consentita vengono accettate. & Origine cross-origin del browser $\leftrightarrow$ CORSMiddleware (FastAPI) \\
		\hline
		\textbf{TI-007} & Verifica$_G$ dell'integrazione$_G$ del pannello dei comandi con la logica di streaming e lo stato globale, il click su ``Genera'' disabilita il bottone tramite isGenerating, un errore del provider compare nell'Alert tramite errorMessage e il testo presente nell'editor resta intatto. & AIPanel $\leftrightarrow$ useLlmStream $\leftrightarrow$ Store Zustand \\
		\hline
	\end{tabularx}
	\caption{Test di integrazione}
	\label{tab:test-integrazione}
\end{table}


\subsection{Test di sistema}

I test$_G$ di sistema$_G$ verificano il comportamento complessivo dell'applicazione rispetto ai Requisiti Funzionali$_G$ e non funzionali. Essi vengono definiti a partire dai casi d'uso$_G$ e dai requisiti$_G$ presenti nell'Analisi dei Requisiti$_G$.

Per il progetto \textit{Second Brain}, i test$_G$ di sistema$_G$ devono coprire almeno i seguenti scenari:
\begin{itemize}
    \item scrittura di testo Markdown nell'area di editing;
    \item visualizzazione corretta del testo formattato;
    \item applicazione della funzione di riassunto;
    \item applicazione della funzione di riscrittura;
    \item applicazione della funzione di traduzione;
    \item applicazione delle critiche secondo i sei cappelli per pensare$_G$;
    \item generazione di un testo completo a partire da prompt;
    \item salvataggio e caricamento di una nota.
\end{itemize}

\subsection{Test di regressione}

I test$_G$ di regressione vengono eseguiti dopo modifiche al codice, correzioni di difetti o introduzione di nuove funzionalità. Il loro obiettivo è verificare che il comportamento già validato non sia stato compromesso.

Nel caso del progetto \textit{Second Brain}, i test$_G$ di regressione sono particolarmente importanti perché le funzionalità basate su LLM$_G$ e prompt possono essere raffinate più volte durante lo sviluppo$_G$. Ogni modifica ai prompt, alle chiamate API$_G$ o alla gestione della risposta deve quindi essere accompagnata dalla riesecuzione dei test$_G$ collegati.

\subsection{Test di accettazione}

I test$_G$ di accettazione verificano che il prodotto rispetti le aspettative del committente$_G$ e soddisfi i requisiti$_G$ obbligatori previsti dal capitolato$_G$.

Essi sono formulati dal punto di vista dell'utente finale e descrivono scenari completi di utilizzo. Il superamento dei test$_G$ di accettazione costituisce una condizione necessaria per considerare il prodotto conforme agli obiettivi minimi del progetto.

\section{Tracciamento dei test}

Il tracciamento dei test$_G$ permette di collegare ogni verifica$_G$ al requisito corrispondente. In questo modo è possibile controllare quali requisiti$_G$ risultano coperti e quali richiedono ulteriori attività di verifica$_G$.

Ogni test$_G$ viene identificato secondo la seguente convenzione:

\begin{itemize}
    \item \textbf{TU}: test$_G$ unitario;
    \item \textbf{TI}: test$_G$ di integrazione$_G$;
    \item \textbf{TS}: test$_G$ di sistema$_G$;
    \item \textbf{TA}: test$_G$ di accettazione.
\end{itemize}

\subsection{Test di sistema principali}

\begingroup
\small
\renewcommand{\arraystretch}{1.25}
\begin{longtable}{|p{1.8cm}|p{9cm}|p{2.5cm}|p{1.5cm}|}
    \hline
    \textbf{ID} & \textbf{Descrizione} & \textbf{Requisito} & \textbf{Stato} \\
    \hline
    \endfirsthead

    \hline
    \textbf{ID} & \textbf{Descrizione} & \textbf{Requisito} & \textbf{Stato} \\
    \hline
    \endhead

    TS-001 & Verificare che l'Utente possa inserire testo libero nell'area di editing in formato Markdown$_G$. & R-1-F-Ob & NI \\
    \hline
    TS-002 & Verificare che l'Utente possa selezionare una porzione di testo presente nell'editor. & R-2-F-Ob & NI \\
    \hline
    TS-003 & Verificare che l'Utente possa tagliare una porzione di testo selezionata, rimuovendola dall'editor e memorizzandola negli appunti di sistema$_G$. & R-3-F-Ob & NI \\
    \hline
    TS-004 & Verificare che l'Utente possa copiare una porzione di testo selezionata negli appunti di sistema$_G$ senza modificare il contenuto dell'editor. & R-4-F-Ob & NI \\
    \hline
    TS-005 & Verificare che l'Utente possa incollare nell'editor un contenuto testuale precedentemente copiato o tagliato. & R-5-F-Ob & NI \\
    \hline
    TS-006 & Verificare che l'Utente possa annullare l'ultima modifica effettuata nell'editor. & R-6-F-Ob & NI \\
    \hline
    TS-007 & Verificare che l'Utente possa ripristinare un'operazione precedentemente annullata, se non sono state effettuate nuove modifiche successive. & R-7-F-Ob & NI \\
    \hline
    TS-008 & Verificare che l'Utente possa applicare e rimuovere la formattazione in grassetto al testo tramite marcatori Markdown$_G$. & R-8-F-Ob & NI \\
    \hline
    TS-009 & Verificare che l'Utente possa applicare e rimuovere la formattazione in corsivo al testo tramite marcatori Markdown$_G$. & R-9-F-Ob & NI \\
    \hline
    TS-010 & Verificare che l'Utente possa applicare e rimuovere la formattazione sottolineata al testo. & R-10-F-Ob & NI \\
    \hline
    TS-011 & Verificare che l'Utente possa applicare e rimuovere la formattazione barrata al testo tramite marcatori Markdown$_G$. & R-11-F-Ob & NI \\
    \hline
    TS-012 & Verificare che l'Utente possa inserire e rimuovere titoli e sottotitoli mediante sintassi Markdown$_G$. & R-12-F-Ob & NI \\
    \hline
    TS-013 & Verificare che l'Utente possa inserire ed eliminare elenchi puntati o numerati. & R-13-F-Ob & NI \\
    \hline
    TS-014 & Verificare che l'Utente possa inserire e rimuovere blocchi di citazione. & R-14-F-Ob & NI \\
    \hline
    TS-015 & Verificare che l'Utente possa inserire ed eliminare blocchi di codice sorgente nel documento. & R-15-F-Ob & NI \\
    \hline
    TS-016 & Verificare che l'Utente possa inserire ed eliminare formule scientifiche o matematiche. & R-16-F-Ob & NI \\
    \hline
    TS-017 & Verificare che l'Utente possa inserire ed eliminare collegamenti ipertestuali. & R-17-F-Ob & NI \\
    \hline
    TS-018 & Verificare che l'Utente possa inserire, modificare ed eliminare tabelle Markdown$_G$. & R-19-F-Ob & NI \\
    \hline
    TS-019 & Verificare che l'Utente possa inserire, sostituire ed eliminare immagini all'interno del documento. & R-20-F-Ob & NI \\
    \hline
    TS-020 & Verificare che l'Utente possa alternare la visualizzazione tra solo editor, solo anteprima e visualizzazione combinata. & R-25-F-Ob & NI \\
    \hline
    TS-021 & Verificare che l'Utente possa ottenere un riassunto del testo selezionato tramite LLM$_G$. & R-27-F-Ob & NI \\
    \hline
    TS-022 & Verificare che l'Utente possa tradurre il testo selezionato in una lingua supportata tramite LLM$_G$. & R-29-F-Ob & NI \\
    \hline
    TS-023 & Verificare che l'Utente possa modificare lo stile linguistico del testo selezionato tramite LLM$_G$. & R-30-F-Ob & NI \\
    \hline
    TS-024 & Verificare che l'Utente possa correggere grammaticalmente il testo selezionato tramite LLM$_G$. & R-31-F-Ob & NI \\
    \hline
    TS-025 & Verificare che l'Utente riceva una comunicazione nel caso in cui l'LLM$_G$ non rilevi errori nel testo analizzato. & R-32-F-Ob & NI \\
    \hline
    TS-026 & Verificare che l'Utente possa analizzare criticamente il testo secondo la metodologia dei Sei Cappelli per Pensare$_G$. & R-33-F-Ob & NI \\
    \hline
    TS-027 & Verificare che l'Utente possa selezionare una specifica prospettiva di analisi tra Cappello Bianco, Rosso, Giallo, Nero, Verde e Blu. & R-34-F-Ob & NI \\
    \hline
    TS-028 & Verificare che l'Utente possa generare automaticamente un testo a partire da un prompt$_G$. & R-35-F-Ob & NI \\
    \hline
    TS-029 & Verificare che l'Utente possa salvare, scartare o rigenerare l'output$_G$ prodotto dall'LLM$_G$. & R-40-F-Ob & NI \\
    \hline
    TS-030 & Verificare che il Sistema$_G$ gestisca gli errori di comunicazione o risposta del servizio LLM$_G$, preservando il contenuto originale dell'editor. & R-42-F-Ob & NI \\
    \hline
    TS-031 & Verificare che il Sistema$_G$ impedisca l'invio di richieste AI quando il testo fornito è insufficiente per un'elaborazione significativa. & R-43-F-Ob & NI \\
    \hline
    TS-032 & Verificare che l'Utente possa creare una nuova nota vuota nell'editor. & R-44-F-Ob & NI \\
    \hline
    TS-033 & Verificare che il Sistema$_G$ gestisca eventuali errori durante la creazione di una nuova nota, informando l'Utente e preservando lo stato precedente. & R-46-F-Ob & NI \\
    \hline
    TS-034 & Verificare che l'Utente possa caricare una nota da file locale in formato .md o .txt. & R-47-F-Ob & NI \\
    \hline
    TS-035 & Verificare che il Sistema$_G$ gestisca eventuali errori durante il caricamento di una nota, preservando la sessione di lavoro precedente. & R-49-F-Ob & NI \\
    \hline
    TS-036 & Verificare che l'Utente possa salvare una nota sul File System locale. & R-50-F-Ob & NI \\
    \hline
    TS-037 & Verificare che il Sistema$_G$ gestisca eventuali errori durante il salvataggio di una nota, preservando l'integrità dei dati presenti nell'editor. & R-52-F-Ob & NI \\
    \hline
    TS-038 & Verificare che l'Utente possa eliminare una nota dal File System locale. & R-53-F-Ob & NI \\
    \hline
    TS-039 & Verificare che il Sistema$_G$ gestisca eventuali errori durante l'eliminazione di una nota, evitando la perdita inconsistente dei dati. & R-55-F-Ob & NI \\
    \hline
    TS-040 & Verificare che l'Utente possa rinominare una nota locale. & R-56-F-Ob & NI \\
    \hline
    TS-041 & Verificare che il Sistema$_G$ gestisca eventuali errori durante la rinominazione di una nota, mantenendo il titolo precedente. & R-58-F-Ob & NI \\
    \hline

    TS-042 & Verificare che l'Utente possa inserire e rimuovere ancore di navigazione interne al documento. & R-18-F-De & NI \\
    \hline
    TS-043 & Verificare che l'Utente possa attivare e disattivare l'autocompletamento dei simboli durante la digitazione. & R-21-F-De & NI \\
    \hline
    TS-044 & Verificare che l'Utente possa attivare e disattivare la sincronizzazione dello scorrimento tra editor e anteprima. & R-22-F-De & NI \\
    \hline
    TS-045 & Verificare che l'Utente possa alternare l'interfaccia tra tema chiaro e tema scuro. & R-23-F-De & NI \\
    \hline
    TS-046 & Verificare che l'Utente possa aprire e chiudere la barra laterale dell'interfaccia. & R-24-F-De & NI \\
    \hline
    TS-047 & Verificare che l'Utente possa inserire testo tramite dettatura vocale, se il browser supporta le Web Speech API$_G$. & R-26-F-De & NI \\
    \hline
    TS-048 & Verificare che l'Utente possa scegliere la percentuale di concisione del riassunto. & R-28-F-De & NI \\
    \hline
    TS-049 & Verificare che l'Utente possa definire la lunghezza desiderata del testo generato. & R-36-F-De & NI \\
    \hline
    TS-050 & Verificare che il Sistema$_G$ mostri all'Utente un riepilogo del contesto prima dell'invio della richiesta all'LLM$_G$. & R-37-F-De & NI \\
    \hline
    TS-051 & Verificare che l'Utente possa fornire un URL come fonte di contesto per la generazione automatica di testo. & R-38-F-De & NI \\
    \hline
    TS-052 & Verificare che l'Utente possa annullare una richiesta di elaborazione in corso. & R-41-F-De & NI \\
    \hline
    TS-053 & Verificare che l'Utente possa visualizzare le informazioni e i metadati di una nota selezionata. & R-59-F-De & NI \\
    \hline

    TS-054 & Verificare che l'Utente autenticato possa selezionare manualmente note salvate nel database$_G$ da usare come contesto per la generazione automatica. & R-39-F-Op & NI \\
    \hline
    TS-055 & Verificare che l'Utente autenticato possa creare una nuova nota nel database$_G$. & R-45-F-Op & NI \\
    \hline
    TS-056 & Verificare che l'Utente autenticato possa caricare una nota dal database$_G$. & R-48-F-Op & NI \\
    \hline
    TS-057 & Verificare che l'Utente autenticato possa salvare una nota nel database$_G$. & R-51-F-Op & NI \\
    \hline
    TS-058 & Verificare che l'Utente autenticato possa eliminare una nota dal database$_G$. & R-54-F-Op & NI \\
    \hline
    TS-059 & Verificare che l'Utente autenticato possa rinominare una nota salvata nel database$_G$. & R-57-F-Op & NI \\
    \hline
    TS-060 & Verificare che l'Utente autenticato possa cercare note nel database$_G$ per titolo. & R-60-F-Op & NI \\
    \hline
    TS-061 & Verificare che l'Utente autenticato possa effettuare una ricerca full-text nel contenuto delle note salvate nel database$_G$. & R-61-F-Op & NI \\
    \hline
    TS-062 & Verificare che l'Utente autenticato possa filtrare le note del database$_G$ per data di creazione. & R-62-F-Op & NI \\
    \hline
    TS-063 & Verificare che l'Utente autenticato possa filtrare le note del database$_G$ per data di ultima modifica. & R-63-F-Op & NI \\
    \hline
    TS-064 & Verificare che l'Utente non autenticato possa effettuare l'autenticazione tramite username e password. & R-64-F-Op & NI \\
    \hline
    TS-065 & Verificare che il Sistema$_G$ gestisca eventuali errori di autenticazione, mantenendo l'Utente disconnesso e mostrando un messaggio di errore. & R-65-F-Op & NI \\
    \hline
    TS-066 & Verificare che l'Utente non autenticato possa registrare un nuovo account. & R-66-F-Op & NI \\
    \hline
    TS-067 & Verificare che il Sistema$_G$ impedisca la registrazione con uno username già presente nel database$_G$. & R-67-F-Op & NI \\
    \hline
    TS-068 & Verificare che il Sistema$_G$ impedisca la registrazione se la password e la conferma password non coincidono. & R-68-F-Op & NI \\
    \hline
    TS-069 & Verificare che l'Utente autenticato possa eliminare il proprio account e i dati cloud associati. & R-69-F-Op & NI \\
    \hline
    TS-070 & Verificare che l'Utente autenticato possa effettuare il logout. & R-70-F-Op & NI \\
    \hline

\end{longtable}

\endgroup

\subsection{Test di accettazione principali}

\begingroup
\small
\renewcommand{\arraystretch}{1.25}
\begin{longtable}{|c|p{7cm}|c|c|}
    \hline
    \textbf{ID} & \textbf{Descrizione} & \textbf{Requisito} & \textbf{Stato} \\
    \hline
    \endhead

    TA-001 & L'utente scrive una nota Markdown, visualizza l'anteprima formattata e salva il contenuto prodotto. & RF-OBB-001, RF-OBB-002, RF-OBB-015 & NI \\ \hline
    TA-002 & L'utente seleziona una porzione di testo e richiede al sistema$_G$ di riassumerla tramite LLM$_G$. & RF-OBB-003, RF-OBB-004 & NI \\ \hline
    TA-003 & L'utente richiede una riscrittura del testo con tono diverso e verifica$_G$ l'inserimento del risultato nell'editor. & RF-OBB-005 & NI \\ \hline
    TA-004 & L'utente traduce una nota in una lingua diversa da quella originale. & RF-OBB-006 & NI \\ \hline
    TA-005 & L'utente applica le analisi dei sei cappelli e riceve una risposta coerente con il punto di vista scelto. & RF-OBB-007 -- RF-OBB-012 & NI \\ \hline
    TA-006 & L'utente fornisce un prompt e ottiene un testo completo generato dal modello linguistico. & RF-OBB-013, RF-OBB-014 & NI \\ \hline
    TA-007 & L'utente carica una nota salvata in precedenza e riprende la modifica del contenuto. & RF-OBB-016 & NI \\ \hline

    \caption{Test di accettazione principali.}
    \label{tab:test-accettazione-principali}
\end{longtable}
\endgroup

\section{Cruscotto di valutazione}

Il cruscotto di valutazione raccoglie periodicamente i valori delle metriche definite nel presente documento. La sua funzione è fornire una visione sintetica dello stato del progetto e rendere più semplice individuare eventuali scostamenti rispetto agli obiettivi fissati.

Il cruscotto viene aggiornato al termine di ogni periodo di lavoro rilevante e viene utilizzato durante le retrospettive per:
\begin{itemize}
    \item valutare l'efficacia delle attività svolte;
    \item individuare problemi ricorrenti;
    \item stabilire eventuali azioni correttive;
    \item migliorare la pianificazione futura;
    \item verificare l'effetto delle iniziative di miglioramento.
\end{itemize}

Le misurazioni che seguono fanno riferimento al periodo che intercorre tra l'aggiudicazione del capitolato, avvenuto il 28 marzo 2026, e la consegna della RTB.

\subsection{Confronto MPC-PV e MPC-EV}
\begin{table}[H]
    \centering
	\begin{tabular}{|c|l|l|l|l|}
		\hline
		\textbf{Sprint} & \textbf{Planned Value} & \textbf{Earned Value} & \textbf{Accumulo PV} & \textbf{Accumulo EV} \\
		\hline
		1 & €985,00 & €939,19 & €985,00 & €939,19 \\ 
		\hline
        2 & €980,00 & €1.024,55 & €1.965,00 & €1.963,73 \\
        \hline
        3 & €955,00 & €1.085,23 & €2.929,00 & €3.048,96 \\
        \hline
        4 & €880,00 & €832,43 & €3.800,00 & €3.881,39 \\
        \hline
        5 & & & & \\
        \hline
	\end{tabular}
	\caption{MPC-PC e MPC-EV per sprint}
	\label{tab:pv-ev}
\end{table}

Come possiamo notare dalla figura §\ref{fig:pv-ev}, il \textit{Planned Value} cresce di pari passo con l'avanzare degli sprint. Vediamo poi come nel primo sprint l'\textit{Earned Value} sia minore rispetto alla stima pianificata: questo è dovuto da una inesperienza nella programmazione del lavoro che, a partire dallo sprint 2 è stata migliorata. 
Questo lo si può notare anche dall'accumulo di EV che, dallo sprint 3, comincia ad essere maggiore dell'accumulo di PV, che riporta quindi un avanzamento del progetto di poco in anticipo rispetto a quanto pianificato. Valutazione:
\begin{itemize}
    \item \textbf{Valore accettabile:} PV$>0$€, vincolo rispettato.
    \item \textbf{Valore ottimo:} PV$<$EV, vincolo rispettato in quasi tutti gli sprint.
\end{itemize}
\begin{figure}[H]
    \includegraphics[width=1\textwidth]{charts/pv-ev.pdf}
    \caption{Comparazione \textit{Planned Value - Earned Value}}
    \label{fig:pv-ev}
\end{figure}

\subsection{Confronto MPC-AC e MPC-ETC}
\begin{table}[H]
    \centering
	\begin{tabular}{|c|l|l|l|p{3cm}|}
		\hline
		\textbf{Sprint} & \textbf{Actual Cost} & \textbf{Accumulo AC} & \textbf{Estimate to Complete} & \textbf{EAC} \\
		\hline
		1 & €945,00 & €945,00 & €11.994,61 & €12.939,61 \\ 
		\hline
        2 & €980,00 & €1.925,0 & €10.375,87 & €12.300,87 \\
        \hline
        3 & €1.105,00 & €3.030,00 & €10.064,31 & €13.094,31 \\
        \hline
        4 & €780,00 & €3.810,00 & €8.239,99 & €12.049,99 \\
        \hline
        5 & & & & \\
        \hline
	\end{tabular}
	\caption{MPC-AC, MPC-ETC e MPC-EAC per sprint}
	\label{tab:ac-etc}
\end{table}
Dal grafico in figura \ref{fig:ac-etc}, possiamo notare come l'\textit{Actual Cost} continui a salire, il che è del tutto normale in quanto man mano che il progetto avanza si consumano risorse economiche. Per quanto riguarda invece l'\textit{Estimate to Complete} e l'\textit{Estimate at Completion}, notiamo che nel primo sprint le curve scendevano, un segnale positivo. Tra il secondo e il terzo sprint, tuttavia, notiamo un peggioramento delle performance di progetto: i costi reali continuano a crescere senza un proporzionale calo del lavoro residuo (ETC), portando la stima del costo finale (EAC) a salire. Questo è successo a seguito della revisione dell'\textit{Analisi dei Requisiti} che ha impegnato buona parte del team senza portare un avantamento effettivo del progetto. A questo si aggiunge una consistente inesperienza nella programmazione del lavoro di gruppo che ha gravato sul bilancio.
Fortunatamente, dal terzo sprint in poi si registra un netto recupero dell'efficienza: l'ETC subisce un calo deciso e, di conseguenza, la previsione del costo totale finale torna ad abbassarsi in modo consistente, il che è un buon segno. Valutazione:
\begin{itemize}
    \item \textbf{Valore accettabile:} AC$>0$€, vincolo rispettato.
    \item \textbf{Valore ottimo:} AC$\leq$EAC, vincolo rispettato.
    \item \textbf{Valore accettabile:} ETC$\geq 0$, vincolo rispettato.
    \item \textbf{Valore ottimo:} ETC$\geq 0$, vincolo rispettato.
\end{itemize}
\begin{figure}[H]
    \includegraphics[width=1\textwidth]{charts/ac-etc.pdf}
    \caption{Comparazione \textit{Actual Cost - Estimate to Complete - Estimated at Completion}}
    \label{fig:ac-etc}
\end{figure}

\end{document}